\section{When to Use Model-Based RL}
% Slide 41
\begin{frame}{When to use model-based RL?}
\begin{itemize}
    \item Advantages:
    \begin{itemize}
        \item Models are immensely useful if easy to learn
        \item Model can be trained without reward labels (self-supervised)
        \item Model is somewhat task-agnostic (can sometimes be transferred
              across rewards)
    \end{itemize}
\end{itemize}
\end{frame}

% Slide 42
\begin{frame}{When to use model-based RL?}
\begin{itemize}
    \item Advantages:
    \begin{itemize}
        \item Models are immensely useful if easy to learn
        \item Model can be trained without reward labels (self-supervised)
        \item Model is somewhat task-agnostic (can sometimes be transferred
              across rewards)
    \end{itemize}
    \item Disadvantages:
    \begin{itemize}
        \item Models don't optimize for task performance
        \item Sometimes harder to learn than a policy
    \end{itemize}
\end{itemize}
\end{frame}

% Slide 43
\begin{frame}{When to use model-based RL?}
\begin{itemize}
    \item Advantages:
    \begin{itemize}
        \item Models are immensely useful if easy to learn
        \item Model can be trained without reward labels (self-supervised)
        \item Model is somewhat task-agnostic (can sometimes be transferred
              across rewards)
    \end{itemize}
    \item Disadvantages:
    \begin{itemize}
        \item Models don't optimize for task performance
        \item Sometimes harder to learn than a policy
    \end{itemize}
\end{itemize}
\vspace{1.5em}
\begin{center}
\textcolor{red}{Whether to use a model depends on how hard it is to learn!}
\end{center}
\end{frame}
